\documentclass{article}
\usepackage[utf8]{inputenc}
\usepackage[dutch]{babel}

\begin{document}

We kennen het verhaal van de blinden en de olifant. De één voelt een slurf en zegt: ``een slang.'' De ander voelt een poot en zegt: ``een boom.'' Een derde voelt een oor en zegt: ``een waaier.'' Iedereen ervaart iets echts, maar niemand overziet het geheel.

De gebruikelijke les is dat individuele waarneming beperkt is. Maar volgens mij zit de werkelijk belangrijke les ergens anders.

Geen mens beschikt over een intern belletje dat afgaat zodra hij ``de waarheid'' heeft gevonden. Dat geldt voor individuen, maar ook voor groepen, instituties en misschien straks voor AI-systemen. Wij ervaren overtuiging, geen garantie op juistheid.

Daardoor wordt de manier waarop we met verschillen omgaan geen morele luxe, maar een praktische noodzaak.

Wat gebeurt er namelijk als de blinden met elkaar gaan praten? Wint de luidste stem? De meerderheid? Alleen wat iedereen bevestigt? Of behandelen ze tegenspraak als informatie dat hun beeld misschien onvolledig is? Zijn ze bereid van positie te wisselen?

Dat verschil bepaalt niet alleen de sfeer van het gesprek. Het bepaalt letterlijk wat een groep in staat is te begrijpen.

Onder ``de sterkste stem heeft gelijk'' ontstaat een denkbeeldig dier dat lijkt op de ervaring van de dominantste persoon. Onder ``alleen consensus telt'' blijft er nauwelijks iets over, omdat niemand hetzelfde deel aanraakt. Maar zodra verschillen niet worden weggewerkt maar onderzocht, ontstaat langzaam een betrouwbaarder beeld van de werkelijkheid.

Dat mechanisme lijkt veel verder te reiken dan een parabel.

Democratie werkt alleen zolang afwijkende signalen niet direct als vijandig worden behandeld. Wetenschap werkt alleen zolang kritiek mag corrigeren. Vrede vraagt dat groepen tijdelijk kunnen verdragen dat hun eigen verhaal misschien niet volledig is. En AI-systemen zullen uiteindelijk niet veilig worden doordat ze simpelweg ``het juiste antwoord'' hebben, maar doordat ze gecorrigeerd kunnen blijven worden.

Misschien volgt ethiek daaruit niet als een verheven ideaal, maar als iets veel fundamentelers: een noodzakelijke strategie voor wezens die nooit rechtstreeks kunnen weten dat ze gelijk hebben.

Eerlijk beschrijven wat je waarneemt. Ruimte laten voor correctie. Tegenspraak behandelen als signaal in plaats van bedreiging. Bereid zijn van perspectief te wisselen zodra nieuwe informatie daarom vraagt.

Niet omdat dat moreel netjes klinkt, maar omdat systemen die dat niet doen uiteindelijk het contact met de werkelijkheid verliezen.

Misschien is dat geen vrijblijvende deugd, maar een van de diepste praktische voorwaarden voor samenleven die er bestaan.

\end{document}